\section{Conclusions}
\label{sec:05_Conclusions}

This work successfully demonstrates an integrated approach to aerodynamic analysis, combining high-fidelity numerical simulation with a data-driven surrogate model. The primary achievement was the extension of the \textbf{\lifex} library to handle turbulent external aerodynamic flows through the integration of the \textbf{Spalart-Allmaras RANS model}. This new capability was validated on the \textbf{Leonardo HPC cluster}, establishing a robust pipeline for generating reference-quality simulation data. Additionally, a \textbf{Convolutional Neural Network} was developed and trained on this data, creating a surrogate model capable of near-instantaneous aerodynamic prediction.\\

The CFD simulations of the NACA airfoil produced lift coefficients that were physically consistent, though moderately underestimated compared to some experimental data. This behavior is physically consistent with the modeling choices made: the VMS-LES formulation employed in the Navier-Stokes step acts as an implicit filter that dissipates small-scale turbulent structures, suppressing the contribution of fine-grained vortical activity to the overall lift generation. Nonetheless, the results remain well within the range of published numerical and experimental values, confirming the physical validity of the simulation framework by \cite{ladson1988} and \cite{mccroskey1988}.\\

A further modeling assumption concerns the compressibility of the flow. The simulations were performed at a Mach number significantly below $0.3$, a regime in which the relative density variations induced by pressure fluctuations are negligible, making the incompressible Navier-Stokes equations an appropriate model.\\

Regarding the deep learning implementation, the results demonstrate that a relatively simple Convolutional Neural Network architecture can guarantee robust performance, achieving a Mean Squared Error of $0.0104$ on the validation set. The primary advantage of this approach lies in its computational efficiency: while a high-fidelity CFD simulation requires approximately 3 hours on an HPC node, the trained CNN inference time is under one second. Even when accounting for the 5 minutes and 11 seconds required for training, this deep learning framework achieves a remarkable 35.11$\times$ speedup compared to conventional methods. This outcome strongly aligns with the premises established by \cite{Cuozzo2026AIDriven}, extending their insights to regular profiles, whereas their approach specifically accounted for deformed airfoils due to ice accretion.\\

This contrast between computational cost and prediction speed places our work within the context of modern aerodynamic design. From a broader perspective, the high-fidelity approach developed here occupies the opposite end of the cost spectrum from purely surrogate-based methods. For instance, the framework proposed by \citet{Cuozzo2026AIDriven} achieves a $10^3-10^4\times$ reduction in wall-clock time at the cost of a $\sim14\%$ underestimation of the iced lift curve slope. The present work, instead, prioritizes physical fidelity to provide the reference-quality data that such surrogate models ultimately depend on for training and validation. The two approaches are therefore complementary rather than competing, and the framework developed here constitutes a natural candidate as a high-fidelity data source for future surrogate-based pipelines targeting iced airfoil configurations.\\

Nevertheless, certain limitations inherent to each method must be acknowledged. These include the computational demand of the RANS simulations and, for the data-driven model, the absence of an \textit{a priori} error estimate, the strict dependency on a training dataset that accurately resembles the target inference domain\footnote{As noted, since the training data from the CFD was slightly biased, we expect this bias persists in the surrogate model.}, and the general lack of interpretability characteristic of black-box models.\\

Looking forward, the framework developed here opens up promising avenues for engineering design and analysis. The speed of the trained CNN makes it ideal for real-time applications and rapid design iterations on consumer-grade hardware, completely bypassing the need for constant HPC access. This capability is particularly relevant for the agile aerodynamic design and field-testing of components for Unmanned Aerial Vehicles and quadcopters, where on-the-fly performance evaluation can accelerate development cycles. The simulation pipeline itself is now poised to serve as a data engine for training future surrogate models on more complex geometries, such as the iced airfoils that motivated this study.


\section{Contributors}
\label{sec:06_Contributors}

\begin{table}[H]
    \centering
    \resizebox{\textwidth}{!}{%
    \begin{tabular}{|l|ccccccccccc|}
        \hline
        \scriptsize\textbf{Name}
        & \rotatebox{90}{\scriptsize \textbf{CNN API}}
        & \rotatebox{90}{\scriptsize \textbf{Topology}}
        & \rotatebox{90}{\scriptsize \textbf{Dropout}}
        & \rotatebox{90}{\scriptsize \textbf{Physics Weights}}
        & \rotatebox{90}{\scriptsize \textbf{Report}}
        & \rotatebox{90}{\scriptsize \textbf{Data Struct.}}
        & \rotatebox{90}{\scriptsize \textbf{Optimizer}}
        & \rotatebox{90}{\scriptsize \textbf{Regularization}}
        & \rotatebox{90}{\scriptsize \textbf{Activation}}
        & \rotatebox{90}{\scriptsize \textbf{Learning Rate}}
        & \rotatebox{90}{\scriptsize \textbf{Training}} \\
        \hline

        \scriptsize\textbf{Giulio Enzo Donninelli}
        & $\checkmark$ & $\checkmark$ & & & $\checkmark$
        & $\checkmark$ & & $\checkmark$ & $\checkmark$ & & $\checkmark$ \\
        \hline

        \scriptsize\textbf{Alessia Rigoni}
        & $\checkmark$ & & $\checkmark$ & $\checkmark$ & $\checkmark$
        & & $\checkmark$ & & & $\checkmark$ & \\
        \hline

        \scriptsize\textbf{Vittorio Sironi}
        & & & $\checkmark$ & $\checkmark$ & &
        & & & & & \\
        \hline

        \scriptsize\textbf{Alessandro Verrengia}
        & & $\checkmark$ & & & &
        & & $\checkmark$ & & & \\
        \hline

    \end{tabular}%
    }
    \caption{Contributions to the work presented in this report. A checkmark indicates participation in the corresponding task.}
    \label{tab:contributors}
\end{table}
